% !TEX root = ../thesis.tex

\chapter{Systémová príručka}\label{ch:system.manual}

Táto kapitola je určená správcom systému a vývojárom, ktorí potrebujú
nainštalovať, nakonfigurovať a spustiť webovú aplikáciu kvantového generátora
náhodných čísel. Príručka predpokladá základnú znalosť príkazového riadku
a prostredia jazyka Python.

% -----------------------------------------------------------------------
\section{Požiadavky na systém}\label{sec:system.requirements}

Pred inštaláciou aplikácie je potrebné mať na cieľovom počítači nainštalované
\textbf{Python~3.10} alebo novší vrátane správcu balíčkov \texttt{pip}.

% -----------------------------------------------------------------------
\section{Štruktúra projektu}\label{sec:project.structure}

Zdrojový kód aplikácie je uložený v adresári
\texttt{bakalarska-praca-main/quantum\_rng}.
Všetky príkazy uvedené v tejto kapitole sa vykonávajú z tohto koreňového
adresára projektu, pokiaľ nie je uvedené inak.

% -----------------------------------------------------------------------
\section{Inicializácia databázy}\label{sec:db.migrate}

Prvým krokom je vytvorenie databázovej schémy pomocou zabudovaného
migračného systému Django. Príkaz \texttt{migrate} aplikuje všetky dostupné
migrácie a vytvorí potrebné tabuľky:

\begin{verbatim}
python manage.py migrate
\end{verbatim}

Po úspešnom dokončení je databáza pripravená na ďalšie kroky.

% -----------------------------------------------------------------------
\section{Vytvorenie správcovského účtu}\label{sec:superuser}

Administrátorský účet sa vytvára interaktívnym príkazom:

\begin{verbatim}
python manage.py createsuperuser
\end{verbatim}

Systém postupne vyzve na zadanie prihlasovacieho mena, e-mailovej adresy
a hesla. Príklad úspešného priebehu:

\begin{verbatim}
Username: admin
Email address: admin@example.com
Password: ********
Password (again): ********
Superuser created successfully.
\end{verbatim}

Takto vytvorený účet má plný prístup k administrátorskému rozhraniu
aplikácie.

% -----------------------------------------------------------------------
\section{Spustenie vývojového servera}\label{sec:runserver}

Po inicializácii databázy a vytvorení správcovského účtu je možné spustiť
zabudovaný vývojový server príkazom:

\begin{verbatim}
python manage.py runserver
\end{verbatim}

Aplikácia je následne dostupná v prehliadači na adrese
\texttt{http://127.0.0.1:8000/}.

% -----------------------------------------------------------------------
\section{Prvotné prihlásenie a nastavenie profilu}\label{sec:first.login}

Po spustení servera je potrebné prihlásiť sa do aplikácie pomocou účtu
vytvoreného v predchádzajúcom kroku. Po prihlásení systém presmeruje
používateľa na stránku profilu (\texttt{/accounts/profile/}), kde je možné
vyplniť základné údaje.

Rola \textbf{admin} sa nastavuje v kolónke \emph{Profil} v sekcii \emph{Admin Editable Fields}.
Používateľom s rolou administrátora je následne sprístupnený odkaz
\emph{Admin Panel} v navigácii aplikácie.

% -----------------------------------------------------------------------
\section{Načítanie vzorových dát}\label{sec:load.data}

Administrátorský panel je dostupný na adrese \texttt{/administration/}.
Nachádza sa v ňom tlačidlo \textbf{Load Local Results}, ktoré načíta
vzorové výsledky uložené z kvantového počítača do lokálnej databázy.
Po úspešnom načítaní sú tieto dáta dostupné vo všetkých sekciách
aplikácie, vrátane štatistických grafov a histórie generovania.

% -----------------------------------------------------------------------

\section{Prehľad URL adries aplikácie}\label{sec:url.overview}

Aplikácia poskytuje nasledujúce URL adresy:

\begin{itemize}
  \item \texttt{/} -- hlavná stránka aplikácie,
  \item \texttt{/coin/} -- generátor náhodného hodu mincou,
  \item \texttt{/dice/} -- generátor náhodného hodu kockou,
  \item \texttt{/generator/} -- všeobecný generátor náhodných čísel
  v zadanom rozsahu,
  \item \texttt{/charts/} -- štatistické grafy a vyhodnotenie generátora,
  \item \texttt{/about/} -- informačná stránka o projekte,
  \item \texttt{/administration/} -- administrátorský panel
  (prístup len pre administrátorov),
  \item \texttt{/accounts/profile/} -- profil prihláseného používateľa,
  \item \texttt{/accounts/history/} -- história generovania náhodných čísel.
\end{itemize}

% -----------------------------------------------------------------------
\section{Postup inštalácie -- zhrnutie}\label{sec:install.summary}

Pre prehľadnosť uvádzame celý postup inštalácie a spustenia v postupnosti
krokov:

\begin{enumerate}
  \item Prejdite do adresára projektu:
  \texttt{cd bakalarska-praca-main/quantum\_rng}
  \item Inicializujte databázu:
  \texttt{python manage.py migrate}
  \item Vytvorte správcovský účet:
  \texttt{python manage.py createsuperuser}
  \item Spustite vývojový server:
  \texttt{python manage.py runserver}
  \item Otvorte prehliadač na adrese \texttt{http://127.0.0.1:8000/}
  a prihláste sa.
  \item Nastavte rolu \emph{admin} v sekcii profilu.
  \item V administrátorskom paneli (\texttt{/administration/}) načítajte
  vzorové dáta tlačidlom \textbf{Load Local Results}.
\end{enumerate}
